% Sources --- the one home of the attributions of section 3 (Q-0428, the fixes round; the register
% review of sections 3.3 to 4, structural change 1, flags 1, 8, 16, 17, 18; the author's direction
% of 2026-10-12: attribution exact and made once). Paper-only: nothing here is shared with the
% blueprint. Until this round the table stood in the opening of section 3.3, before every
% statement it lists, with a third column that graded each statement ("restated", "extended",
% "proved"); the column is gone, and the cells are the source, its anchor and the case that the
% source covers. The statements of the body carry no sentence of attribution, but for two names a
% reader needs to use a statement (Watanabe after the shape; Halgreen for the Student-t clause).
%
% Anchors, each from the note that read it at the page (records/affine/), all confirmed by the
% citation audit of 2026-10-12 (CITATION-AUDIT-paper-iso-2026-10-12.md, 45 sites, all right):
% sato1999levy, Thm. 16.1, p. 99, by Sato 1991, p. 117, and the proof of Thm. 8.1(ii), p. 40
% (READING-sato1999levy-prior-art.md); hazod2001stable, Thm. 2.14.21, p. 383, with its limit
% condition (READING-hazod2001stable.md; the audit's note on row 38); knopova2013note, p. 126 and
% Lemma 2.6, p. 132 (READING-knopova2013note.md); grzywny2019asymptotic, Example 3, Thm. 3.5,
% Prop. 5.3, (5.4), by the numbers of the published article (READING-grzywny2019asymptotic.md);
% sato1978distribution, Lemma 5.1, p. 298, Thm. 1.7, p. 277 (READING-sato1978distribution.md);
% bingham1987regular, Thm. 3.7.3, pp. 162-163; watanabe1983isoperimetric, Proposition, p. 488
% (READING-watanabe1983isoperimetric.md: the law at each time of a symmetric Levy process whose
% Levy measure is isotropic unimodal; the audit's note on row 39); leonenko2013tauberian, Thms. 3
% and 4, arXiv v1 (READING-leonenko2013tauberian.md, with its addendum: the survey is where the
% theorems are collected; the citation is to the book); bingham1997abel, Thms. 1.1 and 1.2 and
% (1.4) of the preprint (READING-bingham1997abel.md, with its addendum: both for a function
% ultimately decreasing to zero); cygan2017asymptotic,
% Thms. 8 and 9, arXiv v2 (READING-cygan2017asymptotic.md); jurek1983integral, Thm. 3.2, p. 252,
% Thm. 4.1, p. 254, Thm. 7.2, p. 259 (READING-jurek1983integral-driving-process.md);
% sato1984operator, Thm. 4.1, p. 86; jurek1992operator, Cor. 1.7, pp. 1056-1057;
% jurek1993operator, Prop. 3.9.6, pp. 180-182; schoenberg1938metric, section 2.1 and Thm. 2,
% pp. 816-817; gneiting1999radial, p. 94; cambanis1981theory, Lemma 2, p. 371, Cor. 2, p. 372
% (the radius of a projection); fang1990symmetric, Thm. 2.10, p. 36 (the density of the
% projection itself, the audit's note on row 17; never its remark of p. 37); the cone paper,
% Lem. 4.12 and Rem. 4.10; trabs2014infinitely, Prop. 1, p. 58.
%
% Q-0430 (the primary sources, read at the page images; records/affine/
% CITATIONS-primary-sources-A.md and -B.md with their reading notes, and
% CITATION-CHANGES-q0430.md for every citation added, changed or left out, with the reason):
% leonenko1999limit, Thm. 1.4.3, p. 65 (the implication from the covariance function to the
% spectral function: 0 < alpha < n, n >= 1, no monotonicity), its proof on p. 104 from
% Lemma 1.4.11, p. 72; Thm. 1.4.4, pp. 66-67 (the book prints the monotonicity of the density
% as "decreasing for ||lambda|| >= k", the survey has (0, eps]: the cell says "a nonincreasing
% function of the radius" and quotes neither; nor the book's remark on c_5(1, alpha), p. 67,
% nor the constant of its Example 1.4.1, both misprinted); the credits of its Comments, p. 351;
% leonenko1991tauberian, Corollary of Lemma 2, p. 1542 of the translation (a pure power, the
% constant L(n) a_n(alpha)/Gamma(1 + alpha) with L(n) "a constant");
% soni1975slowlya, Thms. 1 and 2, p. 169, with the remark of p. 171; soni1975slowlyc,
% Thms. 17(a) and 18(a), pp. 613-614, the Hankel kernel on pp. 617-618 (its display of the
% coefficients c_j is off by a factor sqrt 2: no constant is taken from it). Bingham 1972 and
% Ridenhour and Soni 1974, to which Bingham and Inoue give their Theorem 1.2, are not cited:
% the first is a theorem for another transform, the second the converse implication under
% other conditions (READING-bingham1972tauberian.md, READING-ridenhour1974parseval.md).
% sato1983stationary, Thms. 5.1-5.3, pp. 546-547; jurek1982integral, Thms. 2.1 and 3.1,
% folios 388 and 390; wolfe1982continuous, Thm. 1, p. 306, Cor. 3, p. 307;
% daley2014dimension, Prop. 11 with (3.8), p. 1821, Lemma 12, p. 1822 (its Thm. 4, p. 1815, was
% read in the text layer only and is not cited).
% The versions cited (the author, 2026-10-12: the copies held are the ones the paper cites):
% cygan2017asymptotic and leonenko2013tauberian in the numbering of their arXiv versions (v2,
% v1), bingham1997abel in that of the Hokkaido preprint (no. 370), said once, in the caption.
% Their published numbering was not compared. grzywny2019asymptotic is cited by the statement
% numbers of the published article and never by a page.
%
% Q-0450 (the citation audit of 2026-10-13, CITATION-AUDIT-paper-iso-2026-10-13.md, 54 sites):
% two rows corrected. bogdan2014density, (19): p. 11 of the arXiv version held (v2), not p. 9,
% which has (17) and (18); read again at the page for this item. watanabe1983isoperimetric,
% Proposition, p. 488: its condition (3) is S = aI with a >= 0 and the Levy measure isotropic
% unimodal, so the row and the sentence after Proposition prop:isotropic-radial-monotone name
% the Gaussian covariance as well.

\subsection{Sources}\label{sec:isotropic-sources}

Table~\ref{tab:in-print} says where the main statements of \S\ref{sec:isotropic} are in print,
statement by statement, with the case that each source covers. Several of the sources are
written for a \Levy{} process (\S\ref{sec:line-bridge}): with a
law of exponent $F_d$ they study the laws with the transforms $e^{-\tau F_d}$, $\tau > 0$, the
laws of the process at the times $\tau$. The time is not the scale of this paper. The time
multiplies the exponent, and the scale rescales the frequency. A kernel is the law at time $1$,
and the law at time $\tau$ is a kernel with $\tau c$ in place of $c$.

% Q-0430: a longtable (it is longer than a page and breaks between rows). The struts of a row
% are those of the table's text size: \arraystretch is the ratio of the two line spacings.
{\renewcommand{\arraystretch}{0.81}
\begin{longtable}{@{}S{0.2\textwidth}S{0.76\textwidth}@{}}
\toprule
statement & where it is in print, and for which case \\
\midrule
\endhead
\bottomrule
\endfoot
\bottomrule
\addlinespace
\caption{Where the main statements of \S\ref{sec:isotropic} are in print, and for which case.
Three sources are cited in the numbering of the version that was read, which is not the
journal's: the Hokkaido preprint for \sscite{bingham1997abel}, and the arXiv versions for
\sscite{cygan2017asymptotic} and \sscite{leonenko2013tauberian}. The numbers of
\sscite{grzywny2019asymptotic} are those of the published article, and the pages of
\sscite{leonenko1991tauberian} those of the English translation.}
\label{tab:in-print}
\endlastfoot
Thm.~\ref{thm:isotropic-classification} and Cor.~\ref{cor:isotropic-scale-zero-kernels}, the
classification &
  Read on laws: a process with independent increments on $\RR^d$ that is self-similar up to a
  change of scale and a deterministic shift has self-decomposable laws, and every
  self-decomposable law but a point mass is the law at time $1$ of one
  (\sscite{sato1999levy}, Thm.~16.1, p.~99, due to Sato \sscite{sato1991selfsimilar}). On
  groups, a law is self-decomposable exactly when it is the limit of the kernels $\mu_{0,t}$ of
  a continuous hemigroup that a contracting one-parameter group of automorphisms maps into
  itself (\sscite{hazod2001stable}, Thm.~2.14.21, p.~383). \\
\addlinespace
Thm.~\ref{thm:origin-boundedness-d}, the threshold &
  On the line: \S\ref{sec:line-origin}. For the laws of a \Levy{} process on $\RR^d$: bounded
  when $c > d$, unbounded when $c \le d/2$ (\sscite{knopova2013note}, p.~126, and the proof of
  Lemma~2.6, p.~132). The threshold with its boundary case $c = d$, for two families: the
  \Matern{} kernels (the sources quoted in Example~\ref{ex:matern}, where two of the kernels
  above the threshold are given with their derivatives) and the laws with transform
  $(1 + |\xi|^\alpha)^{-\tau}$ (\sscite{grzywny2019asymptotic}, Example~3). \\
\addlinespace
Cor.~\ref{cor:origin-profile-exponent}, $c$ read off the exponent &
  Clause (1), for every infinitely divisible law: \sscite{sato1999levy}, proof of Thm.~8.1(ii),
  p.~40. Clause (2) for a finite $c$: on the line, an asymptotic form of the transform with its
  constant (\sscite{sato1978distribution}, Lemma~5.1, p.~298); in $\RR^d$ it follows from an
  equivalence for the \Levy{} processes whose laws are rotation invariant with radially
  nonincreasing densities (\sscite{grzywny2019asymptotic}, Thm.~3.5) by de Haan's theorem on
  functions whose increments are governed by a slowly varying function
  (\sscite{bingham1987regular}, Thm.~3.7.3, pp.~162--163). \\
\addlinespace
Prop.~\ref{prop:isotropic-radial-monotone}, the shape &
  For the law at each time of a symmetric \Levy{} process whose Gaussian covariance is a
  multiple of the identity and whose \Levy{} measure has a radially
  nonincreasing density: \sscite{watanabe1983isoperimetric}, Proposition, p.~488. That a
  self-decomposable law on $\RR^d$ which is not degenerate is absolutely continuous:
  \sscite{sato1999levy}, Thm.~27.13, p.~181, proved by Sato in 1982 (p.~196). \\
\addlinespace
Lemma~\ref{prop:origin-rate-d} of \suppl, the rate through the transform: clause (1) for
$c < d$, quoted before Proposition~\ref{prop:origin-rate-profile-d}, and clause (2) for $c = d$ &
  Clause (1), for a rotation-invariant probability law on $\RR^d$, $d \ge 1$, whose transform is
  $|\xi|^{-c}$ times a slowly varying factor, $0 < c < d$: the mass of a ball
  (\sscite{leonenko1999limit}, Thm.~1.4.3, p.~65, where the law is the spectral measure of a
  random field and the transform its covariance function), and the density when it is a
  nonincreasing function of the radius (the same book, Thm.~1.4.4, pp.~66--67). The book
  gives the two theorems for a transform asymptotic to a power to Leonenko and Olenko, and
  their proofs with a slowly varying factor to Olenko (p.~351). For a power the mass of a
  ball is \sscite{leonenko1991tauberian}, Corollary, p.~1542, with a constant that is not given
  in closed form. The survey \sscite{leonenko2013tauberian} collects both theorems, as its
  Thms.~3 and~4, with related ones. \\
\addlinespace
Lemma~\ref{prop:origin-rate-d}, through Hankel transforms &
  For $c > (d-3)/2$ the mass of a ball in clause (1) also follows from
  the Abel--Tauber theorem for Hankel transforms, in the direction from a function that is
  ultimately decreasing to zero to its transform. For a bounded function, which is the case
  $d \ge 3$: \sscite{soni1975slowlya}, Thm.~1, p.~169, with the remark on p.~171, and, as an
  equivalence for a function decreasing on the half-line, \sscite{soni1975slowlyc},
  Thm.~17(a), p.~613, with pp.~617--618 for the Hankel kernel. With local integrability in
  place of boundedness, and so for every $d$: \sscite{bingham1997abel}, Thm.~1.2. The mass of
  a ball in clause (2) follows from the boundary case of that theorem
  (\sscite{soni1975slowlya}, Thm.~2, p.~169, and \sscite{soni1975slowlyc}, Thm.~18(a),
  pp.~613--614, for a bounded function; \sscite{bingham1997abel}, Thm.~1.1 with (1.4)). \\
\addlinespace
Lemma~\ref{prop:origin-rate-d}, on the line, and the method of its proof &
  On the line: the theorem of Sato and Yamazato quoted in \S\ref{sec:line-origin}. The method
  of proof, an average of the transform against a positive kernel read as a Laplace transform of
  the mass of balls: with the Poisson kernel, for the spectral measure of a random field,
  \sscite{leonenko1999limit}, Lemma~1.4.11, p.~72, and the proof on p.~104; with the Gaussian
  kernel, in $\RR^d$, for the measure of the expected time that a \Levy{} process spends in a
  set, \sscite{cygan2017asymptotic}, Thm.~8 with its proof, and Thm.~9. The constant of the
  mass of a ball in clause (1), $R(d,c)\,|S^{d-1}|/c$, is the constant of these two theorems at
  the index $c$. \\
\addlinespace
Prop.~\ref{prop:origin-rate-profile-d}, the rate through the profile &
  On the line: the same theorem of Sato and Yamazato, with their lemma cited at
  Corollary~\ref{cor:origin-profile-exponent} above. In $\RR^d$, for the \Levy{} processes whose
  laws have radially nonincreasing densities and as the time $\tau$ tends to $0$: the density at
  $x$ is asymptotic to $\tau$ times the \Levy{} density at $x$ times the transform of the law at
  the frequencies of length $1/|x|$ (\sscite{grzywny2019asymptotic}, Thm.~3.9). The constant
  $R(d,\tau c)$ of the proposition for the law at time $\tau$ agrees with it in that limit. \\
\addlinespace
Props.~\ref{prop:origin-lower-bound-d} and~\ref{prop:origin-upper-bound-d}, the two bounds;
Lemma~\ref{lem:isotropic-exponent-log} &
  For the \Levy{} processes whose laws have radially nonincreasing densities, a lower bound and the
  comparison of a radially nonincreasing density with its average, on which the upper bound rests
  (\sscite{grzywny2019asymptotic}, Prop.~5.3 and (5.4)). The lemma on the line:
  \sscite{trabs2014infinitely}, Prop.~1, p.~58, and the same lemma of Sato and Yamazato. The
  two-sided estimates of \sscite{bogdan2014density} for such processes assume an exponent whose
  growth lies between two powers of positive order ((19), p.~11, of the arXiv version), which a
  finite $c$ excludes. \\
\addlinespace
Prop.~\ref{prop:isotropic-stages}, the stages &
  On the line, for the \Matern{} kernel of index $1$: the stage between two scales of ratio $r$
  as a mixture in which no displacement has the weight $r^2$ (\prog{B}{Prop.~7.2}, with (7.2)). \\
\addlinespace
Prop.~\ref{prop:radial-dictionary}, the dictionary &
  Clauses (1) and (2) are a case of the representation of a self-decomposable law as the limit
  law of the integrals $\int_0^t e^{-s}\,dY(s)$ of a \Levy{} process $Y$
  (\sscite{jurek1983integral}, Thms.~3.2 and~4.1, pp.~252 and~254; \sscite{sato1984operator},
  Thm.~4.1, p.~86): the case in which $Y$ moves by isolated jumps, at the rate $c$ and with the
  law of $X$. The representation is also in \sscite{sato1983stationary}, Thms.~5.1--5.3,
  pp.~546--547 (the law of a stationary process of Ornstein--Uhlenbeck type), in
  \sscite{jurek1982integral}, Thms.~2.1 and~3.1, pp.~388 and~390 (operator-self-decomposable
  laws on a Banach space), and, on the line, in \sscite{wolfe1982continuous}, Thm.~1, p.~306,
  and Cor.~3, p.~307 (the limit in law of $e^{-\gamma t}\int_0^t e^{\gamma u}\,dB(u)$,
  $\gamma > 0$). In radial form: \sscite{jurek1992operator}, Cor.~1.7, pp.~1056--1057, restated as
  \sscite{jurek1993operator}, Prop.~3.9.6, pp.~180--182. The classes of the radial
  characteristic functions on $\RR^m$ decrease in $m$, with the mixtures of Gaussians as their
  intersection (\sscite{schoenberg1938metric}, \S2.1 and Thm.~2, pp.~816--817), and strictly
  (\sscite{gneiting1999radial}, p.~94). \\
\addlinespace
Prop.~\ref{prop:profile-one-dimension-down} and Cor.~\ref{cor:dictionary-filtration}, the test
on the \Levy{} measure &
  For probability laws, the radius of a projection of a rotation-invariant law
  (\sscite{cambanis1981theory}, Lemma~2, p.~371, and Cor.~2, p.~372) and the density of the
  projection (\sscite{fang1990symmetric}, Thm.~2.10, p.~36). See also
  \sscite{daley2014dimension}: Prop.~11 with (3.8), p.~1821, for the radius law of a radial
  characteristic function in a lower dimension, the law of its radius times an independent
  variable with values in $(0,1)$. The formula of Lemma~\ref{lem:shadow-polar-profile} is (3.8)
  with the \Levy{} measure in place of the law. And Lemma~12, p.~1822: $\Omega_d$ is a
  characteristic function on $\RR^d$ and not on $\RR^{d+1}$, which is the last assertion of
  Corollary~\ref{cor:dictionary-filtration}(2). The representation of a
  self-decomposable law of which the $\Cin$ rays are the elementary terms on the line:
  \sscite{jurek1983integral}, Thm.~7.2, p.~259. \\
\addlinespace
Thm.~\ref{thm:separating-family-critical-values}, the separating family &
  The family, the members $w = 8$ and $w = 4$ as witnesses in dimensions $1$ and $2$, and the
  first critical values, computed numerically: \prog{B}{Lem.~4.12 and Rem.~4.10}. \\
\end{longtable}
}
